% 27_body.tex — 산출물 ㉗ AI 기능 Eval Set 명세서 · 오류 예산 · HITL 승인 설계 (빈 템플릿 / 완성 예시 공용 본문) — gen_product_templates.py 가 Product_specs.py 에서 생성
% 드라이버: 27_Eval명세오류예산HITL_빈템플릿.tex (\wbblanktrue) / 27_Eval명세오류예산HITL_완성예시.tex (\wbblankfalse — 워크북 §X.5 확정 후)
\documentclass[10pt,a4paper]{article}
\usepackage[a4paper,margin=16mm,top=14mm,bottom=20mm]{geometry}
\def\DocNum{27}\def\DocDay{6}\def\DocIdx{2}
\def\DocTitle{AI 기능 Eval Set 명세서 · 오류 예산 · HITL 승인 설계}
\def\DocSub{AI Eval Specification, Error Budget \& Human-in-the-loop}
\def\DocVariant{\B{빈 템플릿}{딥게이지 완성 예시}}
\usepackage{wbstyle}
\def\DocCirc{㉗}   % newunicodechar(㉛~㊵ 폴백) 뒤에 정의해야 활성 문자로 읽힌다
\begin{document}
\docheader
 {\Bg{[회사명 · 제품명]}{딥게이지 · GAUGE Inspect 판정 초안}}
 {\Bg{[v1.x / D+n]}{v1.0 / 2026-07-31 (D+151) · 측정 D+270}}
 {\Bg{[작성자 — CTO·CPO]}{오세진 CTO · 윤하경 CPO\rd{7mm}}}
 {\Bg{[승인자 — CEO]}{강민수\rd{7mm}}}

%% ------------------------------------------------------------------ 1
\secbar{1. 대상 기능과 지표 정의}
\guide{치수 OCR / 외관 결함 분류 / 판정 초안. 지표마다 분자·분모를 적는다(FN = 실제 불량 중 양품 판정 ÷ 실제 불량).}
{\footnotesize
\begin{tabularx}{\linewidth}{|L{30mm}|L{30mm}|Y|Y|}\hline
\hd{기능} & \hd{지표} & \hd{분자} & \hd{분모} \\ \hline
\ifwbblank
\rd{10mm}치수 OCR &  &  &  \\ \hline
\rd{10mm}외관 결함 분류 &  &  &  \\ \hline
\rd{10mm}판정 초안(FN) &  &  &  \\ \hline
\rd{10mm}판정 초안(FP) &  &  &  \\ \hline
\else
치수 OCR & 항목 정확도 & 허용 오차 내로 읽힌 치수 항목 수 & 치수 항목 1,050 \\ \hline
외관 결함 분류 & F1(정밀도·재현율 병기) & TP(결함을 결함으로) & 정밀도: 결함 판정 수 / 재현율: 실제 결함 96 \\ \hline
판정 초안(FN) & 불량 → 양품 오판 비율 & 실제 불량 항목 중 초안 '양품' & 실제 불량 항목 296 \\ \hline
판정 초안(FP) & 양품 → 불량 오판 비율 & 실제 양품 항목 중 초안 '불량' & 실제 양품 항목 1,184 \\ \hline
\fi
\end{tabularx}}

%% ------------------------------------------------------------------ 2
\secbar{2. 데이터셋 구성 — 모집단 정의}
\guide{세트 수·건수·출처·수집 기간·라벨링 규칙·라벨러/검수자. **이 데이터셋이 대표하는 모집단**(어느 산업·어느 공정·어느 조도)을 한 문장으로 적는다 — 이 문장이 Day8까지 따라간다.}
{\footnotesize
\begin{tabularx}{\linewidth}{|L{36mm}|Y|}\hline
\hd{항목} & \hd{내용} \\ \hline
\ifwbblank
\rd{9mm}세트 구성(세트 × 건수) &  \\ \hline
\rd{9mm}출처·수집 기간 &  \\ \hline
\rd{9mm}모집단 정의 — 이 셋이 대표하는 것 &  \\ \hline
\rd{9mm}대표하지 않는 것 &  \\ \hline
\rd{9mm}라벨링 규칙 &  \\ \hline
\rd{9mm}라벨러 / 검수자 &  \\ \hline
\rd{9mm}갱신 주기 &  \\ \hline
\else
세트 구성(세트 × 건수) & 3세트 × 20건 = 60건(V-07) — A 세영정공(절삭) · B 동보프레스(프레스) · C 한빛금속(단조). 1,480항목(치수 1,050 · 외관 430) \\ \hline
출처·수집 기간 & 파일럿 3사 실제 라인, D+160~200, 검사원 스마트폰 촬영, 오전 9~11시 \\ \hline
모집단 정의 — 이 셋이 대표하는 것 & 자동차 2·3차 협력사의 절삭·프레스·단조 공정에서, 오전 자연광+형광등 조도로, 검사원이 스마트폰으로 찍은 초·중·종품 사진 \\ \hline
대표하지 않는 것 & 오후 3시 이후·야간 조도(X-02: 84.7 / 88.2\%) · 사출·주조·조립 · 자동차 외 산업 · 고정 카메라·산업용 조명 · 마스터가 없는 현장 \\ \hline
라벨링 규칙 & 라벨러 2인 독립 라벨 → 불일치 시 강민수 확정(불일치율 6.8\%). 기준서 = 각 사 검사기준서 최신판(B-05 매핑본) \\ \hline
라벨러 / 검수자 & 라벨러 2인(D+190 합류, 전 검사원) / 강민수 / 도구 CTO 스크립트(B-25) \\ \hline
갱신 주기 & 분기 1회 + 모델 갱신 시 + 사고 시. DS-v2 후보: 오후 조도 세트 20건 추가 \\ \hline
\fi
\end{tabularx}}

%% ------------------------------------------------------------------ 3
\secbar{3. 합격 기준과 개정 절차}
\guide{v1 기준(FN ≤ 8.0\% \& OCR ≥ 90\%)과 개정 절차(누가·언제·무엇을 근거로). 사고 후 기준을 바꾸려면 이 절차를 통과해야 한다(criteria drift 방지).}
{\footnotesize
\begin{tabularx}{\linewidth}{|L{36mm}|Y|}\hline
\hd{항목} & \hd{내용} \\ \hline
\ifwbblank
\rd{9mm}합격 기준 v1 &  \\ \hline
\rd{9mm}측정 시점·주기 &  \\ \hline
\rd{9mm}개정 승인권자 &  \\ \hline
\rd{9mm}개정 근거 요건 &  \\ \hline
\rd{9mm}개정 이력 기록 위치 &  \\ \hline
\else
합격 기준 v1 & FN ≤ 8.0\% 그리고 OCR ≥ 90.0\% — DS-v1 기준(V-06). 참고 지표(기준 아님): FP · F1 · 지연 · 단가 \\ \hline
측정 시점·주기 & 첫 측정 D+270(GA 30일 전) · 분기 1회 · 모델 갱신 시 갱신 전후 비교 \\ \hline
개정 승인권자 & 낮추는 개정(FN 상향·OCR 하향·DS 축소): 대표 승인 + 이사회 보고 + 근거 문서 · 높이는 개정: CTO 전결 \\ \hline
개정 근거 요건 & 낮추는 개정은 '측정 오류의 입증' 또는 '모집단 정의의 변경'만. 사고·일정·영업 압력은 근거가 아니다. DS를 바꾸면 기준 버전도 \\ \hline
개정 이력 기록 위치 & 항목 7 + 판정 로그 헤더에 기준 버전 스탬프 \\ \hline
측정 기록 D+270 & OCR 960/1,050 = 91.4\% ✓ · FN 24/296 = 8.1\%(0.1\%p 초과) · FP 173/1,184 = 14.6\% · F1 0.78 · 1.8초 · 32원 → 불합격. 기준 유지. GA는 '합격 전 베타' \\ \hline
\fi
\end{tabularx}}

%% ------------------------------------------------------------------ 4
\landscapeon
\secbar{4. 채점 루브릭 4차원 × 4수준}
\guide{정확 / 근거 제시 / 일관성 / 안전(오판정의 방향) — 각 수준을 관찰 가능한 문장으로.}
{\footnotesize
\begin{tabularx}{\linewidth}{|L{26mm}|Y|Y|Y|Y|}\hline
\hd{차원} & \hd{1 (불합격)} & \hd{2} & \hd{3} & \hd{4 (모범)} \\ \hline
\ifwbblank
\rd{13mm}정확 &  &  &  &  \\ \hline
\rd{13mm}근거 제시 &  &  &  &  \\ \hline
\rd{13mm}일관성 &  &  &  &  \\ \hline
\rd{13mm}안전 &  &  &  &  \\ \hline
\else
정확 & 판정 또는 치수값이 틀렸다 & 판정은 맞으나 치수값 허용 오차 밖 & 판정·치수 모두 허용 오차 내 & 3 + 결함 유형까지 일치 \\ \hline
근거 제시 & 근거 영역 표시 없음 & 영역이 결함 위치와 다르다 & 영역이 결함을 포함 & 3 + 기준서 항목 번호 연결 \\ \hline
일관성 & 같은 사진 3회 판정이 2종 이상 & 3회 중 1회 다름(신뢰도 표시) & 3회 일치, 신뢰도 편차 > 0.1 & 3회 일치, 편차 ≤ 0.1 \\ \hline
안전 & 불량을 양품으로(FN) & 양품을 불량으로(FP) — 검사원이 잡는다 & 오판 없음, 신뢰도 < 0.85로 '보류' & 오판 없음, 신뢰도 ≥ 0.85 \\ \hline
\fi
\end{tabularx}}
\landscapeoff

%% ------------------------------------------------------------------ 5
\secbar{5. 오류 예산 정책}
\guide{월 허용 FN·FP 건수(또는 비율), 소진 시 조치(자동판정 OFF·릴리스 동결·재학습), 신뢰도 임계값과 그 근거.}
{\footnotesize
\begin{tabularx}{\linewidth}{|L{36mm}|Y|L{50mm}|}\hline
\hd{항목} & \hd{값·규칙} & \hd{근거} \\ \hline
\ifwbblank
\rd{9mm}월 FN 예산 &  &  \\ \hline
\rd{9mm}월 FP 예산 &  &  \\ \hline
\rd{9mm}신뢰도 임계값 &  &  \\ \hline
\rd{9mm}소진 시 조치 &  &  \\ \hline
\rd{9mm}측정·보고 주기 &  &  \\ \hline
\else
월 FN 예산 & 라인별 — 초안 '양품' 중 검사원이 '불량'으로 뒤집은 비율 ≤ 8.0\% & 합격 기준과 같은 숫자, 분모는 다르다. 검사원도 놓친 FN은 여기 없다 → 클레임 \\ \hline
월 FP 예산 & 초안 '불량' 중 검사원이 '양품'으로 뒤집은 비율 ≤ 15\% & v1 FP 14.6\%. FP가 늘면 초안을 무시한다 \\ \hline
신뢰도 임계값 & 0.85 미만은 초안 미표시('판정 보류 — 직접 판정') & DS-v1에서 0.85 미만 구간이 전체 FN의 71\% \\ \hline
소진 시 조치 & (자동) 라인 초안 OFF → 원인 분석 48h → 재학습 → 재측정 합격 후 ON · 초과 기간 릴리스 동결 & 사람이 결정하면 늦다(PM ⑰ 형식) \\ \hline
측정·보고 주기 & 주간(내부) · 월간(고객 품질팀장 — 라인별 요약) & 팀장이 보는 숫자가 있어야 한다 \\ \hline
\fi
\end{tabularx}}

%% ------------------------------------------------------------------ 6
\secbar{6. Human-in-the-loop 승인 설계}
\guide{자동판정 ON 조건 / 검사원 승인 흐름 / 에스컬레이션 / **검사원 책임 한계 문구**(계약·화면에 쓰는 문장). 촉탁직 검사원이 이 흐름을 어떻게 볼지 한 줄.}
{\footnotesize
\begin{tabularx}{\linewidth}{|L{36mm}|Y|}\hline
\hd{항목} & \hd{설계} \\ \hline
\ifwbblank
\rd{10mm}자동판정 ON/OFF 조건 &  \\ \hline
\rd{10mm}검사원 승인 흐름(화면 순서) &  \\ \hline
\rd{10mm}에스컬레이션(누구에게·언제) &  \\ \hline
\rd{10mm}검사원 책임 한계 문구 &  \\ \hline
\rd{10mm}현장 수용성 검토 &  \\ \hline
\else
자동판정 ON/OFF 조건 & 자동판정 기능 없음(B-17 Won't). 초안 표시 조건: 베타 라인(세영 오전 1) 그리고 신뢰도 ≥ 0.85 그리고 오류 예산 미소진 \\ \hline
검사원 승인 흐름(화면 순서) & ① 촬영 → ② 초안 카드(값·양/불·신뢰도·근거 영역) → ③ [승인] 1탭 / [수정] + 사유 1탭 → ④ 로그(B-12) → ⑤ 부적합이면 NCR 패키지 \\ \hline
에스컬레이션(누구에게·언제) & 불일치 연속 3건 → 품질팀장(당일) · 예산 소진 → CTO(즉시, 자동 OFF) · 클레임 → 대표 + 팀장 48시간 \\ \hline
검사원 책임 한계 문구 & 화면: '판정 초안은 참고용입니다. 최종 판정은 검사원이 합니다.' / 계약(㉘ 항목 6): '판정 초안의 오류로 인한 당사 책임은 연 계약액의 10\%(162만 원)를 상한으로 한다.' 같은 화면에 두지 않는다 \\ \hline
현장 수용성 검토 & [세영정공] 박선재(촉탁직)의 눈: '최종 판정은 검사원' = '책임은 너'. 예측: 재촬영·별도 종이 기록. A-05b 미해결. 대응: 수정 사유 1탭 · 수정 건수를 모델 품질 지표로 표기 \\ \hline
\fi
\end{tabularx}}

%% ------------------------------------------------------------------ 7
\secbar{7. 개정 이력}
\guide{버전 / 일자 / 무엇을 바꿨나 / 근거 / 승인자. 빈 표로 시작한다.}
{\footnotesize
\begin{tabularx}{\linewidth}{|L{16mm}|L{22mm}|Y|Y|L{24mm}|}\hline
\hd{버전} & \hd{일자} & \hd{변경 내용} & \hd{근거} & \hd{승인자} \\ \hline
\ifwbblank
\rd{9mm} &  &  &  &  \\ \hline
\rd{9mm} &  &  &  &  \\ \hline
\rd{9mm} &  &  &  &  \\ \hline
\rd{9mm} &  &  &  &  \\ \hline
\else
v1.0 & D+151 & 최초 — FN ≤ 8.0 \& OCR ≥ 90 on DS-v1 & ㉔ X-02 · ㉓ 4순위 & 강민수 \\ \hline
(측정) & D+270 & DS-v1 측정 — FN 8.1 불합격, 기준 유지, '합격 전 베타' & 항목 3 절차 & 오세진 · 윤하경 \\ \hline
(후보) & — & DS-v2: 오후 조도 세트 20건 추가 → 기준 v2 & 대표하지 않는 것 1번 & — \\ \hline
 &  &  &  &  \\ \hline
\fi
\end{tabularx}}

\end{document}